% !TEX root = ../main.tex

% Chapter 2 — Methodik

\chapter{Methodik}
\label{ch:methodik}

\section{Projektmanagement und Projektplan}
\label{sec:pm}

% Rollen, Verantwortlichkeiten, Arbeitspakete, Zeitplan (ca. 1 Seite).

\section{Datengrundlage und Beschaffung}
\label{sec:datengrundlage}

% Kurzbeschreibung des SHP, Datenzugang via FORS, verwendete Files
% (SHPLONG_P_USER, SHPLONG_H_USER, SHP_MP, SHP_SO), Storage und Ordnerstruktur.

\section{Datenkatalog}
\label{sec:datenkatalog}

% Tabellarische Übersicht der Zielvariablen: Name, Kurzbeschreibung,
% Quelldatei, Wertebereich. Vollständige Version im Anhang.

\section{Datenaufbereitung}
\label{sec:aufbereitung}

% Pipeline: Einlesen der SPSS-Dateien, Merging (Person/Haushalt/Master),
% Filtering (Erwachsene, Jahre mit PP04), Variablenkonstruktion
% (BIRTHYGEN, HASKID, ...). Übersicht der R-Skripte.

\section{Datenqualität und Plausibilität}
\label{sec:qualitaet}

% Recodierung der SHP-Missing-Codes zu NA, Fehlendanteile,
% Plausibilitätsprüfungen (Wertebereiche, Verteilung über Wellen).
